% !PRINTED 596
\weberpage{596}

%<<CONT
reduce to the unit group. Hence the case must at all events
occur at some point that the group $P$ is reduced by the adjunction
of a root of a pure equation.

The investigation of this question is extraordinarily simplified
if one transforms it somewhat further\,\footnote{This form of the question was first pointed out by C. \gesperrt{Jordan} (Trait\'e des substitutions p. 386).}.

We have seen in \S.\,162 that the pure equations belong to the
\gesperrt{Abel}ian ones, that all \gesperrt{Abel}ian equations are solvable through a
chain of cyclic equations of prime degree, and these latter
by radicals.

We therefore replace the question of solvability by radicals
by the equivalent one of solvability by roots of \gesperrt{Abel}ian
equations. If we replace the \gesperrt{Abel}ian equation by a chain of
cyclic equations of prime degree, then the first reduction of the
group will occur through the adjunction of the root of one such
equation, and so we stand first of all before the question:

\gesperrt{Under what conditions is the group $P$
of an equation of degree $n$, $f(x) = 0$, reduced by the adjunction
of a root of a cyclic equation of prime degree $m$?}

We do not restrict ourselves here to irreducible equations
$f(x) = 0$, but discuss the question in general, always under the
self-evident assumption that $f(x)$ has no multiple
roots.

It is conceivable that it may be necessary first to adjoin to
the field $\Omega$ various roots of cyclic equations (e.\,g. roots of unity)
which do not alter the group $P$. But since we are now
concerned only with ascertaining the necessary properties of the
group $P$, we assume that all the necessary preparations
have been made, and the field $\Omega$ so constituted
that by the adjunction of a root of the cyclic equation
$\varphi(x) = 0$, of prime degree $m$, the group $P$ reduces to
one of its divisors, $Q$.

If we denote the roots of $\varphi(x) = 0$ by

\[
\varepsilon,\ \varepsilon_1,\ \varepsilon_2 \ldots \varepsilon_{m-1},
\]
then all these quantities are expressible rationally (in $\Omega$)
through an arbitrary one among them, and if $P$ is the group of
$f(x) = 0$
%CONT>>
